\section{Genuine original Fubini--Study families from total-space Laurent sections}

\begin{definition}\label{def:native-tc-original-weighted-fs-data}
\lean{EnvelopeProperty.NativeTestConfiguration.TestConfiguration.weightedFSData}
\leanok
Let $M,L$ be the original invertible modules of a test configuration,
with the original action and fibre-one identification. For $n>0$, a
finite original family $G\subset H^0(\mathcal X,M^{\otimes n})$ and
$r>0$, form the actual original Fubini--Study datum of $L$ with
degree and denominator both $n$, and terms
\[
 (a,i\log r),\qquad (a,i)\in T_G.
\]
Here $T_G$ is the original whole Laurent coefficient family chosen
before all points and radii. This is a datum of the actual original
line $L$, with rational class coefficient one and finite real weights.
\uses{thm:native-tc-global-laurent-fixed-family}
\end{definition}
\begin{proof}
The finite image of the actual coefficient/exponent family consists
of entire original sections of $L^{\otimes n}$ and real weights.
The positive denominator $n$ and degree $n$ satisfy the degree
equation $n=n\cdot1$. No section realization or norm identity is an
input to this construction.
\end{proof}

\begin{theorem}\label{thm:native-tc-original-weighted-fs-evaluation}
\lean{EnvelopeProperty.NativeTestConfiguration.TestConfiguration.originalPowerSectionNorm_eq_globalSectionNorm,EnvelopeProperty.NativeTestConfiguration.TestConfiguration.weightedFSData_eval}
\leanok
At every actual original native point $x$, this entire datum has value
\[
 \frac1n\log\max_{(a,i)\in T_G}
       \{\lVert a(x)\rVert_{L^{\otimes n}}r^i\}.
\]
The equality holds in the extended reals, including empty families,
zero sections, common zeroes and negative Laurent exponents.
\uses{def:native-tc-original-weighted-fs-data,def:native-tc-weighted-laurent-section-norm}
\end{theorem}
\begin{proof}
The native pullback-unit value of an original whole section equals
the fibre section of its actual original section morphism. Hence
the coefficient norms used above are precisely the genuine global
section norms used by the original Fubini--Study datum. Its terms
are the finite image of $T_G$, so the datum's supremum is the
supremum of the original pair values. The logarithm is monotone
and takes zero to negative infinity, so it preserves the finite
supremum, including its empty value. Multiplicativity gives
$\log(\lVert a(x)\rVert r^i)=\log\lVert a(x)\rVert+i\log r$;
$r^i$ is strictly positive for every integer $i$. This proves the
identity without excluding any zero coefficient norm.
\end{proof}

\begin{theorem}\label{thm:native-tc-original-weighted-fs-total-log}
\lean{EnvelopeProperty.NativeTestConfiguration.TestConfiguration.exists_weightedFSData_totalLog_factor}
\leanok
For every fixed positive degree, actual point and $0<r\leq1$, one
real number $A_{n,x,r}$, chosen before every finite original family,
satisfies
\[
 \frac1n\log\max_{s\in G}\lVert s\rVert_{\mathcal X,r,x}
 =A_{n,x,r}+\operatorname{FS}_{G,n,r}(x)
\]
for every finite $G$. The additive comparison is derived from the
actual proper-model norm and actual Laurent coefficients.
\uses{thm:native-tc-weighted-laurent-true-section-norms,cor:native-tc-weighted-laurent-finite-section-norms,thm:native-tc-original-weighted-fs-evaluation}
\end{theorem}
\begin{proof}
Choose the genuine positive factor $c_{n,x,r}$ for all original
families. Set $A_{n,x,r}=n^{-1}\log c_{n,x,r}$, a finite real number.
Apply the extended logarithm to the already proved exact finite-family
norm identity. The logarithm of the positive factor is its real
logarithm; multiplication by the positive inverse degree distributes
over its sum with the weighted maximum logarithm. Substitute the
preceding actual Fubini--Study evaluation formula. The argument also
applies when either maximum is zero.
\end{proof}

\begin{theorem}\label{thm:native-tc-original-weighted-fs-no-common-zero}
\lean{EnvelopeProperty.NativeTestConfiguration.TestConfiguration.exists_eventual_weightedFSData_noCommonZero}
\leanok
If the original fibre-one line is ample, then for all sufficiently
large positive $n$ one can choose a finite original total-space
section family $G_n$, before all original native points, such that
the actual weighted Fubini--Study datum above has no common zero
on the entire original native point space.
\uses{thm:native-tc-original-weighted-fs-evaluation,cor:native-tc-weighted-laurent-finite-section-norms,thm:native-test-configuration-actual-gauss-finite-norm-family}
\end{theorem}
\begin{proof}
Original ampleness and eventual global generation supply one finite
actual total-space family in each sufficiently large degree, whose
true Gauss section-norm maximum is strictly positive at every original
point. At each point, the genuine positive comparison factor times
the original coefficient weighted maximum equals this true maximum.
Thus the coefficient weighted maximum is nonzero. If every original
coefficient section had zero norm there, every weighted term and their
finite supremum would be zero, a contradiction. A term with nonzero
true coefficient norm gives a term with the same section in the actual
weighted datum, proving the no-common-zero assertion simultaneously
on the entire original space.
\end{proof}

\begin{corollary}\label{cor:native-tc-original-rational-fs-families}
\lean{EnvelopeProperty.NativeTestConfiguration.TestConfiguration.weightedFSRadius,EnvelopeProperty.NativeTestConfiguration.TestConfiguration.weightedFSRadius_properties,EnvelopeProperty.NativeTestConfiguration.TestConfiguration.weightedFSData_isRational,EnvelopeProperty.NativeTestConfiguration.TestConfiguration.exists_eventual_weightedFSData_isFinFS}
\leanok
At the fixed original radius $r=e^{-1}$, the actual datum has rational
weights for every original finite family. If $L$ is ample, then in
every sufficiently large positive degree the preceding original family
produces a genuine finite rational Fubini--Study function of the
actual original line $L$ and class coefficient one, on the entire
original native point space.
\uses{def:native-tc-original-weighted-fs-data,thm:native-tc-original-weighted-fs-no-common-zero}
\end{corollary}
\begin{proof}
The fixed radius is positive and at most one, and its real logarithm
is $-1$. Hence the weight of every actual original coefficient with
integer Laurent exponent $i$ is $-i$, the image of a rational number.
Choose the preceding actual eventual family, before all original
points. Its weighted datum has positive denominator, rational weights
and no common native zero. These are exactly the defining properties
of a finite rational Fubini--Study function of the original $L$.
Its degree equation has coefficient one by construction, and its
evaluation is the proved genuine whole original function.
\end{proof}

\begin{remark}
Compatibility of $A_{n,x,r}$ across all degrees remains a separate step.
No degree-independent metric or final envelope identity is asserted here.
\end{remark}
